\pagestyle{myheadings} \setcounter{page}{1} \setcounter{footnote}{0}

\section{~设置模式时间步长} \label{app:tstep}
\newcounters
\vssub

模式时间步长按网格逐一设置，并作为模式定义文件 {\file mod\_def.ww3} 中
模式设置的一部分。这意味着在多网格模式设置中（使用模式驱动程序
{\file ww3\_multi}），每个网格都有自己的时间步长设置。本节给出单个网格
以及镶嵌方法中多个网格的时间步长设置建议。实际网格的实用时间步长设置
示例可参见测试案例 {\file mww3\_case\_01} 到 {\file mww3\_case\_03} 中
使用的各个网格。

\subsection{~单个网格}

基本波浪模式网格需要定义四个时间步长，如本手册第
\pageref{dt_list} 页 \para\ref{sec:basic_num} 所述。通常首先要考虑的是
空间传播的 \cfl\ 时间步长，也就是所考虑网格在 {\file ww3\_grid.inp}
中定义的四个时间步长中的第二个。用于判断数值格式稳定性的临界 \cfl\
数 $C_c$ 定义为[对比式~(\ref{eq:quick_4})]

\begin{equation}
C_c = \frac{c_{g,\max} \Delta t}{\min(\Delta x , \Delta y)} \:\:\: ,
\label{eq:CFLmax}
\end{equation}

\noindent
其中 $c_{g,\max}$ 为最大群速，$\Delta t$、$\Delta x$ 和 $\Delta y$
分别为时间和空间增量。最大群速是最低离散模式频率对应的群速。注意到在
给定频率下，最大群速出现在中间水深，因此该最大速度约为最低离散谱频率
深水群速的 1.15 倍。注意，\cfl\ 数在形式上包含流的影响
[式~(\ref{eq:x_dot})] 和网格移动的影响 [式~(\ref{eq:bal_move})]。
后两种影响在模式内部通过根据流速和网格移动速度动态调整相应的最小时间
步长来处理。因此，用户可以在定义这一最小传播时间步长时忽略流和网格移动。
对这里使用的格式，临界 \cfl\ 数为 1。

第二个要考虑的时间步长是总体时间步长（{\file ww3\_grid.inp} 中标识的
第一个时间步长）。为获得最大的数值精度，该时间步长应设置为小于或等于
上述 \cfl\ 时间步长。然而，特别是在球面网格中，临界 \cfl\ 条件通常只
发生在少数网格点。在大多数网格点中，\cfl\ 数会小得多。在这类网格中，
如果将总体时间步长取为临界 \cfl\ 时间步长的 2 到 4 倍，精度通常不会
明显受损。这样的设置通常会显著改善模式经济性。数值精度的关键在于如何
解释 \cfl\ 数。该数表示信息在单个时间步内传播距离的归一化值。如果信息
在源项施加之前已经传播过多个网格箱，就会产生不精确性。当 \cfl\
$\approx 1$ 且总体时间步长为 \cfl\ 时间步长的四倍时，信息会在源项施加
之前传播过四个网格箱。这可能导致模式不精确。然而，如果最大 \cfl\ 数为
1，但平均 \cfl\ 数只有 0.25，正如许多球面网格中即便对最低频率也会出现
的情况，则信息在一个总体时间步中只传播过一个网格箱，不会产生精度问题。

有效的总体时间步长还应考虑模式强迫数据可用的时间间隔，以及请求模式输出
的时间间隔。如果输入和输出时间步长都是总体时间步长的整数倍，则存在一种
平衡且一致的数值积分方案，尽管模式本身并不要求这一点。在这一考虑中最
重要的是结果的可重复性。如果修改输入或输出时间步长，使其不再是总体模式
时间步长的整数倍，则模式中实际的离散时间推进将被这些输入和输出时间步长
改变，因此可以预期对实际模式结果产生影响。这种影响可能可察觉，但通常
非常小。

第三个要考虑的时间步长是最大折射（以及波数偏移）时间步长。为获得最大的
模式经济性，该时间步长应设置为等于（或大于）总体时间步长。然而，这会使
连续模式时间步中的空间计算和折射计算顺序交替，在强折射情况下可能导致
波浪参数出现周期为 $2 \Delta t$ 的小幅振荡。将最大折射时间步长设置为
总体时间步长的一半，可完全避免这种振荡。考虑到模式中折射项的代价很小，
这通常对模式经济性的影响可以忽略。因此，推荐的折射时间步长为总体模式
时间步长的一半。

设置这一时间步长时需要特别注意。为确保数值稳定性，特征折射速度按
式~(\ref{eq:theta_filter}) 进行滤波。该滤波会在海底地形快速变化的情况
下抑制折射。减小折射时间步长可降低这种滤波的影响。因此，稳妥的做法是
用小得多的谱内模式时间步长测试模式网格，以评估该滤波的影响。

最后要设置的时间步长是 \para\ref{sub:source} 中动力源项积分的最小时间
步长。这是一个安全阀，用来避免源项积分中出现过小而难以承受的时间步长。
根据网格增量大小，该值通常设为 5 到 15s。注意，增大这一时间步长并不
一定改善模式经济性；更大的最小源项积分时间步长会增加积分中的谱噪声，
而这反过来可能会{\em 减小}平均源项积分时间步长！


\subsection{~网格镶嵌}

对于使用 {\file ww3\_multi} 构成镶嵌模式的各个单独网格，时间步长设置的
考虑原则上与上一节讨论的单个网格相同。附加考虑如下：

\begin{itemize}

\item 对于镶嵌方法的管理算法能否正常工作，各个网格的总体时间步长不需要
以任何方式“匹配”。不过，如果相同等级的网格共享总体时间步长，并且不同
等级网格之间的时间步长采用整数比，则会更容易跟踪和预测管理算法的运行，

\item 如果两个相同等级的网格发生重叠，则所需重叠区宽度由数值格式的模板
宽度，以及该格式对最长波分量被调用的次数（总体时间步长与最大 \cfl\
时间步长之比）共同决定。因此，增大总体模式时间步长会改善单个网格的模式
经济性，但同时会增加相同等级网格所需的重叠宽度，从而降低镶嵌方法的
经济性。

\end{itemize}

\bpage \pagestyle{empty}
